\section{Discussion}
\label{sec:discussion}

The headline finding is that the SVKIE multi-prior framework lifts
\textsf{total}-field F1 substantially over a single-prior (lexical-
only) baseline across three upstream architectures and two receipt
corpora.  The lift is not attributable to any single component:
the per-component ablation (Tab.~\ref{tab:ablation_neurips}) shows
each prior contributes a distinct, complementary increment.

The most important architectural property is the verifier's
keyword-independence.  Identity $\mathrm{I}_3$ — the structural
subset-sum identity — fires on every receipt satisfying the
itemisation hypothesis regardless of which summary keywords the
printer chose to emit.  This generalises the keyword-anchored
verifiers used in commercial systems (Microsoft Form Recogniser,
AWS Textract) which fail silently when OCR corrupts
\textsf{SUBTOTAL} or \textsf{CASH}.

The architecture-agnostic claim is empirically supported by the
wrapper-$\Delta$ matrix (Tab.~\ref{tab:wrapper_delta}).  The
verifier produces measurable F1 lift when applied to DONUT,
LayoutLMv3, and the in-house FOCUS-T head, demonstrating that the
framework is not specific to a single backbone.  This is the
property that makes SVKIE a methodology contribution rather than a
single-system optimisation.

The boundary at which the framework is silent is the receipt-
construction-hypothesis violation set: bundling discounts, rounding
adjustments, implicit service charges (Sec.~\ref{sec:svkie_theory}).
On these receipts $\mathrm{I}_3$ correctly abstains and the
framework gracefully degrades to the lexical-prior baseline.  We
view graceful abstention as a desirable property of a verifier:
silently committing the wrong value would be worse than no
verification at all.
